\chapter{Quantisation of the electromagnetic field}
\label{ch:em-quantisation}

Part~I quantised mechanical degrees of freedom; Part~II coupled them
to reservoirs. To describe light --- the photons that drive qubits,
carry readout signals, and constitute the dark-matter candidates of
Part~VI --- we must quantise the electromagnetic field itself. The
route is the one already rehearsed for the harmonic oscillator in
\cref{sec:qho}: expand the field in cavity modes, recognise that
each mode is a harmonic oscillator, and promote its amplitude to a
pair of ladder operators. The quanta of excitation are photons, the
ground state is not empty but seethes with vacuum fluctuations, and
the field amplitude becomes an operator with an irreducible noise
floor. This chapter sets up that machinery: the mode expansion and
canonical quantisation, the Fock states and the electric-field
operator, the quadratures and their vacuum uncertainty, and the
photon-number statistics that distinguish one state of light from
another. It is the foundation for the coherent and squeezed states
of \cref{ch:states-of-light} and the atom--field coupling of
\cref{ch:jaynes-cummings}.

% --------------------------------------------------------------
\section{From classical modes to quantum oscillators}
\label{sec:field-quantisation}

\begin{phenomenon}
The electromagnetic field in a cavity does not hold arbitrary
amounts of energy at a given frequency: it holds energy in discrete
quanta of $\hbar\omega$, the photons. The evidence is a century old
--- the blackbody spectrum, the photoelectric effect, the
photon-counting statistics of a photodetector --- and it points to
the same conclusion: a single mode of the field behaves exactly
like a quantum harmonic oscillator, with a ladder of equally spaced
energy levels and a non-zero ground-state energy.
\end{phenomenon}

\paragraph{Mode expansion of the classical field.}
In a cavity of volume $V$ with perfectly conducting walls, the
source-free vector potential (in Coulomb gauge, $\nabla\cdot\bm A=0$)
can be expanded in the discrete set of normal modes $\bm u_{\bm k,s}(\bm r)$
labelled by wavevector $\bm k$ and polarisation $s$,
\begin{equation}
\label{eq:mode-expansion}
\bm A(\bm r,t) \;=\; \sum_{\bm k,s}\Bigl[
\alpha_{\bm k,s}(t)\,\bm u_{\bm k,s}(\bm r) + \text{c.c.}\Bigr],
\end{equation}
where each mode amplitude obeys $\ddot\alpha_{\bm k,s}=-\omega_k^2\,\alpha_{\bm k,s}$
with $\omega_k=c|\bm k|$ --- the equation of a harmonic oscillator.
The total field energy, integrated over the cavity, is a sum of
independent oscillator energies, one per mode.

\begin{statement}[\textbf{Field modes as quantum oscillators}]
\label{stmt:field-oscillators}
Quantising each mode by the canonical prescription of
\cref{sec:qho} --- promoting the mode amplitude to ladder
operators $\hat a_{\bm k,s},\hat a_{\bm k,s}^\dagger$ with
\begin{equation}
\label{eq:field-commutator}
\comm{\hat a_{\bm k,s}}{\hat a_{\bm k',s'}^\dagger} \;=\; \delta_{\bm k\bm k'}\,\delta_{ss'},
\qquad
\comm{\hat a_{\bm k,s}}{\hat a_{\bm k',s'}} \;=\; 0,
\end{equation}
--- yields the field Hamiltonian
\begin{equation}
\label{eq:field-hamiltonian}
\op H_{\rm field} \;=\; \sum_{\bm k,s}\hbar\omega_k
\Bigl(\hat a_{\bm k,s}^\dagger\hat a_{\bm k,s} + \tfrac12\Bigr).
\end{equation}
Each mode is an independent harmonic oscillator (\cref{stmt:qho-spectrum});
the quantum $\hbar\omega_k$ of excitation in mode $(\bm k,s)$ is a
\emph{photon}, and $\hat a_{\bm k,s}^\dagger$ creates one.
\end{statement}

\paragraph{Derivation.}
The classical mode energy is $E_{\bm k,s}=\tfrac12(\,p_{\bm k,s}^2+\omega_k^2 q_{\bm k,s}^2)$
in terms of canonical position and momentum quadratures
$q_{\bm k,s}\propto\mathrm{Re}\,\alpha_{\bm k,s}$,
$p_{\bm k,s}\propto\mathrm{Im}\,\alpha_{\bm k,s}$. This is identical
to the mechanical oscillator of \cref{sec:qho}, so the same
canonical quantisation applies: impose $\comm{\hat q}{\hat p}=i\hbar$,
define $\hat a=(\omega\hat q+i\hat p)/\sqrt{2\hbar\omega}$, and the
energy becomes \cref{eq:field-hamiltonian}. The only new feature is
that there is one oscillator per mode and the modes are spatial
field patterns rather than the position of a particle.

% --------------------------------------------------------------
\section{Fock states and the field operators}
\label{sec:fock-states}

\begin{statement}[\textbf{Fock (number) states of a single mode}]
\label{stmt:fock}
For a single mode (dropping the $\bm k,s$ labels), the
photon-number operator $\op N=\hat a^\dagger\hat a$ has eigenstates
$\ket n$, the \emph{Fock} or \emph{number} states, with
$\op N\ket n=n\ket n$, $n=0,1,2,\dots$. They are built from the
vacuum $\ket0$ by
\begin{equation}
\label{eq:fock-ladder}
\hat a^\dagger\ket n=\sqrt{n+1}\,\ket{n+1},\quad
\hat a\ket n=\sqrt n\,\ket{n-1},\quad
\ket n=\frac{(\hat a^\dagger)^n}{\sqrt{n!}}\,\ket0,
\end{equation}
form a complete orthonormal basis of the single-mode Hilbert space,
and diagonalise the Hamiltonian
$\op H=\hbar\omega(\op N+\tfrac12)$ with energy
$E_n=\hbar\omega(n+\tfrac12)$.
\end{statement}

\paragraph{The electric-field operator.}
Re-expressing the mode expansion \cref{eq:mode-expansion} in terms
of the ladder operators, the electric field of a single mode (with
mode function normalised over the cavity) becomes the operator
\begin{equation}
\label{eq:efield-operator}
\hat{\bm E}(\bm r) \;=\; \bm{\mathcal E}_0\,\bigl[\hat a\,\bm u(\bm r) + \hat a^\dagger\,\bm u^*(\bm r)\bigr],
\qquad
\mathcal E_0=\sqrt{\frac{\hbar\omega}{2\varepsilon_0 V}},
\end{equation}
where $\mathcal E_0$ is the \emph{electric field per photon} --- the
amplitude of the vacuum field fluctuation in the mode volume $V$. It
is this single-photon field that sets the strength of every
light--matter coupling in \cref{ch:jaynes-cummings}: smaller mode
volume means larger $\mathcal E_0$ and stronger coupling, the design
principle behind cavity and circuit QED.

\paragraph{Vacuum fluctuations and zero-point energy.}
The ground state $\ket0$ has $\op N\ket0=0$ --- no photons --- yet
the field is not classically zero. The energy is
$E_0=\tfrac12\hbar\omega$ per mode, and the field operator has
non-zero variance in the vacuum: $\langle0|\hat{\bm E}|0\rangle=0$
but $\langle0|\hat{\bm E}^2|0\rangle=\mathcal E_0^2\neq0$. These
\emph{vacuum fluctuations} are physical: they drive spontaneous
emission (the Purcell effect of \cref{sec:purcell}), they set the
standard quantum limit on amplifiers (\cref{ch:input-output}), and
their mode-volume dependence is what a small cavity exploits.

\paragraph{Transition.}
The Fock states diagonalise the energy, but they are not the states
the field is usually found in --- a laser or a microwave drive
produces something much closer to a classical wave. To compare
states of light we need the language of \emph{quadratures}, which
turns the abstract ladder operators into measurable field
amplitudes.

% --------------------------------------------------------------
\section{Quadratures and vacuum noise}
\label{sec:quadratures}

\begin{statement}[\textbf{Field quadratures and their uncertainty}]
\label{stmt:quadratures}
Define the dimensionless \emph{quadrature operators}
\begin{equation}
\label{eq:quadratures}
\hat X \;=\; \tfrac12(\hat a + \hat a^\dagger),
\qquad
\hat P \;=\; \tfrac{1}{2i}(\hat a - \hat a^\dagger),
\end{equation}
the in-phase and out-of-phase components of the field relative to a
reference oscillation. They satisfy $\comm{\hat X}{\hat P}=\tfrac{i}{2}$
and hence the uncertainty relation
\begin{equation}
\label{eq:quadrature-uncertainty}
\Delta X\,\Delta P \;\geq\; \tfrac14.
\end{equation}
In the vacuum (and in any coherent state) the noise is equal in the
two quadratures and saturates the bound,
$\Delta X=\Delta P=\tfrac12$ --- the \emph{standard quantum limit}
of field detection.
\end{statement}

\paragraph{The phase-space picture.}
A single mode is naturally drawn in the $X$--$P$ plane. The vacuum
is a unit-area ``ball'' of noise centred at the origin, radius
$\tfrac12$ in each quadrature. A classical field is a displaced
ball; a Fock state is a ring; a squeezed state, as we will see in
\cref{ch:states-of-light}, is an ellipse with one quadrature
squeezed below $\tfrac12$ at the expense of the other. The quadrature
variances are exactly what a homodyne detector
(\cref{sec:cascaded}) measures, which is why this representation is
the working language of continuous-variable quantum optics.

% --------------------------------------------------------------
\section{Photon-number statistics}
\label{sec:photon-statistics}

\begin{phenomenon}
Shine a field on an ideal photon counter and record the number of
clicks in a fixed window. The distribution $P(n)$ of counts is a
fingerprint of the state of light: a number state gives exactly $n$
every time (zero variance), a thermal field gives a broad,
super-Poissonian spread, and --- as the next chapter shows --- a
laser gives the Poissonian distribution in between. The mean and
variance of $P(n)$ are directly measurable and distinguish
classical from non-classical light.
\end{phenomenon}

\begin{statement}[\textbf{Number statistics and the Mandel parameter}]
\label{stmt:mandel}
For a state with photon-number distribution $P(n)=|\langle n|\psi\rangle|^2$,
the mean is $\bar n=\langle\op N\rangle$ and the variance
$(\Delta n)^2=\langle\op N^2\rangle-\bar n^2$. The \emph{Mandel
$Q$ parameter}
\begin{equation}
\label{eq:mandel-Q}
Q \;=\; \frac{(\Delta n)^2 - \bar n}{\bar n}
\end{equation}
classifies the statistics: $Q=0$ for a Poissonian (coherent) field,
$Q>0$ for super-Poissonian (e.g.\ thermal, $Q=\bar n$), and
$-1\leq Q<0$ for sub-Poissonian (non-classical) light, with $Q=-1$
for a Fock state ($\Delta n=0$). Sub-Poissonian statistics have no
classical-field description and are a signature of genuinely quantum
light.
\end{statement}

\paragraph{Thermal light.}
A single mode in thermal equilibrium at temperature $T$ has the
Bose--Einstein distribution
$P(n)=\bar n^n/(1+\bar n)^{n+1}$ with mean
$\bar n=(e^{\hbar\omega/k_BT}-1)^{-1}$ and variance
$(\Delta n)^2=\bar n+\bar n^2$, hence $Q=\bar n>0$:
thermal light is super-Poissonian, with number fluctuations larger
than its mean. This is the photon statistics of the blackbody
background that floods a cryostat and the dark-count budget of the
photon counters in \cref{ch:microwave-absorption}.

% --------------------------------------------------------------
This chapter quantised the electromagnetic field by recognising
each cavity mode as a harmonic oscillator: the mode amplitudes
became ladder operators, the quanta became photons, and the field
acquired a Hamiltonian \cref{eq:field-hamiltonian} with a non-zero
ground-state energy. The Fock states diagonalise the photon number,
the electric-field operator \cref{eq:efield-operator} carries the
single-photon amplitude $\mathcal E_0$ that sets every light--matter
coupling, the quadratures expose the vacuum noise floor, and the
Mandel parameter classifies a state's photon statistics.

With the field quantised, \cref{ch:states-of-light} builds the
states of light that actually occur in the laboratory --- the
coherent states of a laser or microwave drive, and the squeezed
states that beat the standard quantum limit --- and
\cref{ch:jaynes-cummings} couples a single mode to a qubit.

\clearpage
\begin{chaptersummary}

\paragraph{Field quantisation.}
Expanding the cavity field in normal modes makes each mode a
harmonic oscillator. Promoting the amplitudes to ladder operators
with $\comm{\hat a_{\bm k,s}}{\hat a_{\bm k',s'}^\dagger}=\delta_{\bm k\bm k'}\delta_{ss'}$
gives $\op H_{\rm field}=\sum_{\bm k,s}\hbar\omega_k(\hat a_{\bm k,s}^\dagger\hat a_{\bm k,s}+\tfrac12)$.
The quanta are photons.

\paragraph{Fock states and field operator.}
Number states $\ket n$ satisfy $\hat a^\dagger\ket n=\sqrt{n+1}\ket{n+1}$
and diagonalise $\op N=\hat a^\dagger\hat a$. The single-mode field
operator is $\hat{\bm E}=\mathcal E_0[\hat a\bm u+\hat a^\dagger\bm u^*]$
with $\mathcal E_0=\sqrt{\hbar\omega/2\varepsilon_0 V}$ the field per
photon; smaller $V$ means stronger coupling.

\paragraph{Quadratures.}
$\hat X=\tfrac12(\hat a+\hat a^\dagger)$,
$\hat P=\tfrac1{2i}(\hat a-\hat a^\dagger)$ obey
$\comm{\hat X}{\hat P}=i/2$, so $\Delta X\,\Delta P\geq\tfrac14$;
vacuum and coherent states saturate it with
$\Delta X=\Delta P=\tfrac12$ (the standard quantum limit). The
vacuum has zero mean field but non-zero variance --- the fluctuations
that drive spontaneous emission.

\paragraph{Photon statistics.}
The Mandel parameter $Q=[(\Delta n)^2-\bar n]/\bar n$ is $0$ for
Poissonian (coherent), $>0$ for super-Poissonian (thermal,
$Q=\bar n$), and $<0$ for sub-Poissonian (non-classical, $Q=-1$ for
a Fock state).

\end{chaptersummary}

% --------------------------------------------------------------
\section*{Exercises}
\addcontentsline{toc}{section}{Exercises}

\begin{exercise}[Single-mode field energy]
Starting from the classical mode energy
$E=\tfrac12(p^2+\omega^2 q^2)$ and the substitution
$\hat a=(\omega\hat q+i\hat p)/\sqrt{2\hbar\omega}$, derive
$\op H=\hbar\omega(\hat a^\dagger\hat a+\tfrac12)$.
\begin{solution}
Invert: $\hat q=\sqrt{\hbar/2\omega}(\hat a+\hat a^\dagger)$,
$\hat p=-i\sqrt{\hbar\omega/2}(\hat a-\hat a^\dagger)$. Then
$\hat p^2=-\tfrac{\hbar\omega}2(\hat a-\hat a^\dagger)^2$,
$\omega^2\hat q^2=\tfrac{\hbar\omega}2(\hat a+\hat a^\dagger)^2$.
Adding, the $\hat a^2$ and $\hat a^{\dagger2}$ terms cancel and
$\op H=\tfrac{\hbar\omega}2(\hat a\hat a^\dagger+\hat a^\dagger\hat a)
=\hbar\omega(\hat a^\dagger\hat a+\tfrac12)$ using
$\comm{\hat a}{\hat a^\dagger}=1$.
\end{solution}
\end{exercise}

\begin{exercise}[Zero-point energy is unobservable as a shift]
Explain why the $\tfrac12\hbar\omega$ per mode can usually be dropped
from $\op H_{\rm field}$, and one situation where its
mode-dependence is physical.
\begin{solution}
A constant added to the Hamiltonian shifts all energies equally and
does not affect dynamics or transition frequencies, so the
$\sum_{\bm k,s}\tfrac12\hbar\omega_k$ can be subtracted by redefining
the energy zero. Its \emph{dependence on boundary conditions} is
physical, however: changing the cavity geometry changes which modes
exist and hence the total zero-point energy, producing the
measurable Casimir force between conductors.
\end{solution}
\end{exercise}

\begin{exercise}[Field per photon scaling]
Compute $\mathcal E_0=\sqrt{\hbar\omega/2\varepsilon_0 V}$ for a
microwave mode at $\omega/2\pi=5\,$GHz in a $V=1\,$mm$^3$ volume,
and comment on why circuit QED achieves strong coupling.
\begin{solution}
$\hbar\omega=h\cdot5\times10^9\approx3.3\times10^{-24}\,$J;
$\varepsilon_0 V=8.85\times10^{-12}\cdot10^{-9}=8.85\times10^{-21}$.
$\mathcal E_0=\sqrt{3.3\times10^{-24}/(2\cdot8.85\times10^{-21})}
\approx\sqrt{1.9\times10^{-4}}\approx1.4\times10^{-2}\,$V/m.
Circuit QED confines the mode to a quasi-1D transmission line with
effective mode volume far below $\lambda^3$, boosting $\mathcal E_0$
and the dipole coupling $g\propto\mathcal E_0$ into the strong-coupling
regime where $g\gg\kappa,\gamma$.
\end{solution}
\end{exercise}

\begin{exercise}[Vacuum field variance]
Show that $\langle0|\hat E|0\rangle=0$ but
$\langle0|\hat E^2|0\rangle=\mathcal E_0^2$ for a single mode (take
$\bm u$ real and evaluate at an antinode).
\begin{solution}
With $\hat E=\mathcal E_0(\hat a+\hat a^\dagger)$ at the antinode,
$\langle0|\hat E|0\rangle=\mathcal E_0\langle0|(\hat a+\hat a^\dagger)|0\rangle=0$
since $\hat a\ket0=0$ and $\langle0|\hat a^\dagger=0$. For the
square, $\hat E^2=\mathcal E_0^2(\hat a^2+\hat a^{\dagger2}+\hat a\hat a^\dagger+\hat a^\dagger\hat a)$;
only $\hat a\hat a^\dagger=\hat a^\dagger\hat a+1$ has non-zero vacuum
expectation, giving $\langle0|\hat E^2|0\rangle=\mathcal E_0^2$. The
field fluctuates even with no photons present.
\end{solution}
\end{exercise}

\begin{exercise}[Quadrature commutator and uncertainty]
From $\comm{\hat a}{\hat a^\dagger}=1$ derive
$\comm{\hat X}{\hat P}=i/2$ and hence
$\Delta X\,\Delta P\geq\tfrac14$.
\begin{solution}
$\comm{\hat X}{\hat P}=\tfrac{1}{2}\cdot\tfrac{1}{2i}\comm{\hat a+\hat a^\dagger}{\hat a-\hat a^\dagger}
=\tfrac1{4i}(\comm{\hat a}{-\hat a^\dagger}+\comm{\hat a^\dagger}{\hat a})
=\tfrac1{4i}(-1-1)=\tfrac{i}{2}$.
The Robertson relation $\Delta X\,\Delta P\geq\tfrac12|\langle\comm{\hat X}{\hat P}\rangle|=\tfrac12\cdot\tfrac12=\tfrac14$.
\end{solution}
\end{exercise}

\begin{exercise}[Vacuum saturates the bound]
Show $\Delta X=\Delta P=\tfrac12$ in the vacuum, confirming it is a
minimum-uncertainty state.
\begin{solution}
$\langle\hat X\rangle=\langle\hat P\rangle=0$ as in the field-mean
exercise. $\langle\hat X^2\rangle=\tfrac14\langle0|(\hat a+\hat a^\dagger)^2|0\rangle=\tfrac14\langle0|\hat a\hat a^\dagger|0\rangle=\tfrac14$,
so $\Delta X=\tfrac12$; identically $\Delta P=\tfrac12$. Then
$\Delta X\,\Delta P=\tfrac14$, saturating \cref{eq:quadrature-uncertainty}.
\end{solution}
\end{exercise}

\begin{exercise}[Number-state quadrature noise]
Compute $\langle\hat X\rangle$ and $(\Delta X)^2$ in a Fock state
$\ket n$. Why does a number state have no well-defined phase?
\begin{solution}
$\langle n|\hat X|n\rangle=\tfrac12\langle n|(\hat a+\hat a^\dagger)|n\rangle=0$
(off-diagonal). $\langle n|\hat X^2|n\rangle=\tfrac14\langle n|(\hat a\hat a^\dagger+\hat a^\dagger\hat a)|n\rangle=\tfrac14(2n+1)$,
so $(\Delta X)^2=\tfrac14(2n+1)$, growing with $n$, and identically
for $P$. The mean field is zero in every quadrature regardless of
$n$: a Fock state is a ring in phase space with completely random
phase --- number and phase are conjugate, and a sharp photon number
means a maximally uncertain phase.
\end{solution}
\end{exercise}

\begin{exercise}[Mandel $Q$ of a Fock state]
Evaluate the Mandel parameter \cref{eq:mandel-Q} for $\ket n$ and
interpret the sign.
\begin{solution}
A Fock state has $\bar n=n$ and $(\Delta n)^2=0$, so
$Q=(0-n)/n=-1$. The maximally negative value marks the most
strongly sub-Poissonian (non-classical) state: the photon number is
perfectly definite. No classical stochastic field can produce
$Q<0$.
\end{solution}
\end{exercise}

\begin{exercise}[Thermal photon statistics]
For the Bose--Einstein distribution
$P(n)=\bar n^n/(1+\bar n)^{n+1}$, verify $\langle\op N\rangle=\bar n$
and $(\Delta n)^2=\bar n+\bar n^2$, hence $Q=\bar n$.
\begin{solution}
Writing $x=\bar n/(1+\bar n)$, $P(n)=(1-x)x^n$, a geometric
distribution. Its mean is $\langle n\rangle=x/(1-x)=\bar n$ and
variance $\mathrm{Var}=x/(1-x)^2=\bar n(1+\bar n)=\bar n+\bar n^2$.
Then $Q=[(\bar n+\bar n^2)-\bar n]/\bar n=\bar n>0$: thermal light is
super-Poissonian (photon bunching), with fluctuations exceeding the
mean.
\end{solution}
\end{exercise}

\begin{exercise}[Number of thermal microwave photons]
Compute the mean thermal occupancy $\bar n$ of a $5\,$GHz mode at
$T=20\,$mK and at $T=4\,$K, and comment on the implication for
microwave single-photon detection.
\begin{solution}
$\hbar\omega/k_B\approx0.24\,$K. At $20\,$mK,
$\bar n=(e^{0.24/0.02}-1)^{-1}=(e^{12}-1)^{-1}\approx6\times10^{-6}$:
the mode is essentially in vacuum, so a single signal photon is
detectable above the thermal background. At $4\,$K,
$\bar n=(e^{0.06}-1)^{-1}\approx16$: a large thermal background that
would swamp single-photon detection. This is why microwave photon
counters (\cref{ch:microwave-absorption}) must operate at tens of
mK.
\end{solution}
\end{exercise}
